\documentclass[journal_revised_MPC]{subfiles}

\begin{document}

We consider the unconstrained optimization problem
\begin{equation}
    \underset{x \in \bbR^{n}}{\mathrm{minimize}} \quad f(x),
    \label{prob:main_problem}
\end{equation}
where $f\colon \mathbb{R}^n \to \mathbb{R}$ is a continuously differentiable nonconvex function, and $x \in \mathbb{R}^n$ is the optimization variable.
We focus on settings where objective function evaluations are contaminated by bounded and non-diminishing numerical noise~\citep{sunTrustRegionMethod2023}.
Specifically, only inexact function values $\overline{f}(x)$ are available for optimization.
Such inaccuracies are unavoidable in many realistic scenarios, including finite-precision arithmetic, simulation-based evaluations, and stochastic approximations~\citep{shiNoiseTolerantQuasiNewtonAlgorithm2022,moreEstimatingComputationalNoise2011,caflischMonteCarloQuasiMonte1998}.
When function values are contaminated by numerical noise, standard optimization algorithms may lose their theoretical guarantees, behave erratically, or terminate prematurely with unreliable solutions.

In the absence of noise and with accurate function evaluations, quasi-Newton methods, in particular the Broyden--Fletcher--Goldfarb--Shanno (BFGS) method and its Limited-memory variant (L-BFGS), are among the most efficient and widely used algorithms for solving problem~\eqref{prob:main_problem}~\citep{liuLimitedMemoryBFGS1989,nocedal1999numerical}. % to LLM: ignore this git difference.
Their success is largely attributed to fast local convergence and scalability.
Standard implementations~\citep{2020SciPy-NMeth,Okazaki2007libLBFGS} rely on line search procedures that enforce the Wolfe conditions to guarantee sufficient descent and the positive definiteness of the Hessian approximation.

In the presence of numerical noise, these line search conditions may become unreliable.
Differences in function values or directional derivatives can be dominated by numerical errors, leading to unstable step sizes, ill-conditioned Hessian approximations, or abnormal termination~\citep{shiNoiseTolerantQuasiNewtonAlgorithm2022,paquetteStochasticLineSearch2020}.
From a computational perspective, improving the robustness of quasi-Newton methods in noisy environments is therefore crucial.
In particular, it is desirable to design algorithms that (i) retain computational efficiency when function evaluations are accurate, (ii) remain stable when function values are unreliable, and (iii) guarantee convergence to first-order stationary points under mild assumptions.

In this work, we propose a noise-tolerant regularized quasi-Newton method that addresses these challenges.
Our approach combines an Armijo line search, quadratic regularization \citep{uedaRegularizedNewtonMethod2014,kanzowRegularizationLimitedMemory2023}, and an Objective-Function-Free Optimization (OFFO) inspired strategy~\citep{Gratton08022024,grattonOFFOMinimizationAlgorithms2023}. % to LLM: ignore this git difference.
When a sufficient decrease in function value is observed, the method behaves similarly to standard quasi-Newton methods with line search, preserving its practical efficiency.
When numerical noise dominates, we introduce regularization on the approximate Hessian to ensure numerical stability, which yields convergence properties similar to those of OFFO methods, particularly AdaGrad-Norm~\citep{wardAdaGradStepsizesSharp2020}.
The resulting algorithm is particularly well-suited for low-precision or noisy environments.

From a theoretical standpoint, we establish that the proposed method achieves the global convergence rate of $\order{1/\varepsilon^2}$ for smooth nonconvex optimization problems, \MPCchanged{which matches the optimal worst-case order for smooth nonconvex optimization with Lipschitz-continuous gradients~\citep{carmonLowerBoundsFinding2020}}, even when objective function values are corrupted by noise.
Extensive numerical experiments on the CUTEst benchmark collection, across various floating-point precisions and artificial noise settings, \MPCchanged{show that the proposed method remains stable under function-value noise and performs favorably against baselines.}

\subsection{\MPCchanged{Comparison With Related Work}}

\begin{table*}[t]
    \begin{MPCchange}
    \ifdefined\MPCRevision\captionsetup{font={color=red}}\fi
    \centering
    \caption{
        Theoretical comparison of the proposed method and its Newton-type, quasi-Newton, and OFFO competitors.
        OFFO is included as the first-order source of the fallback scaling.
        A \ding{51} denotes theoretical coverage and ``--'' its absence.
        The ``Noisy $f$'' and ``Noisy $g$'' columns indicate whether the corresponding errors are allowed.
        The asterisk indicates that gradient noise is outside our analysis and is considered only in joint-noise stress tests.
        Still, unlike a Wolfe line search, the proposed sufficient-decrease test does not require an exact gradient evaluation at each trial point.
    }
    \label{tab:related_work_comparison}
    \setlength{\tabcolsep}{3pt}
        \begin{tabular}{@{}>{\centering\arraybackslash}p{0.20\linewidth}cccc>{\centering\arraybackslash}p{0.35\linewidth}@{}}
            \toprule
            Algorithm
             & Uses $f$
             & Noisy $f$
             & Noisy $g$
             & Nonconvex
             & Guarantee \\
            \midrule
            \makecell[c]{\textbf{\texttt{Line}} / \textbf{\texttt{SciPy}} \\[-0.4em]\scalebox{0.8}{\citep{liuLimitedMemoryBFGS1989,Okazaki2007libLBFGS,2020SciPy-NMeth}}}
             & \ding{51}
             & --
             & --
             & \ding{51}
             & R-linear under uniform convexity \\
            \makecell[c]{Regularized Newton\\[-0.4em]\scalebox{0.8}{\citep{uedaConvergencePropertiesRegularized2010}}}
             & \ding{51}
             & --
             & --
             & \ding{51}
             & \(\order{\varepsilon^{-2}}\) to \(\norm{\nabla f}\leq\varepsilon\) \\
            \makecell[c]{\textbf{\texttt{Reg}}\\[-0.4em]\scalebox{0.8}{\citep{kanzowRegularizationLimitedMemory2023}}}
             & \ding{51}
             & --
             & --
             & \ding{51}
             & Global stationarity \\
            \makecell[c]{\textbf{\texttt{NTQN}}\\[-0.4em]\scalebox{0.8}{\citep{xieAnalysisBFGSMethod2020,shiNoiseTolerantQuasiNewtonAlgorithm2022}}}
             & \ding{51}
             & \ding{51}
             & \ding{51}
             & --
             & R-linear to the noise floor \\
            \makecell[c]{\textbf{\texttt{OFFO}}\\[-0.4em]\scalebox{0.8}{\citep{Gratton08022024}}}
             & --
             & \ding{51}
             & --
             & \ding{51}
             & \(\order{\varepsilon^{-2}}\) to \(\|\nabla f\|\leq\varepsilon\) \\
            \makecell[c]{\textbf{\texttt{Ours}}\\[-0.4em]\scalebox{0.8}{\cref{alg:regularized_qn}}}
             & \ding{51}
             & \ding{51}
             & --\rlap{\raisebox{0.3ex}{\large *}}
             & \ding{51}
             & \(\order{\varepsilon^{-2}}\) to \(\|\nabla f\|\leq\varepsilon\) \\
            \bottomrule
        \end{tabular}
    \end{MPCchange}
\end{table*}

\begin{MPCchange}
We next compare our method with related quasi-Newton and OFFO methods, whose theoretical guarantees are summarized in \cref{tab:related_work_comparison}.

The closest methods are regularized Newton and quasi-Newton methods~\citep{uedaConvergencePropertiesRegularized2010,kanzowRegularizationLimitedMemory2023}.
These methods assume reliable function values, while we analyze the behavior of our method under persistent bounded relative-or-absolute function errors and give the order-optimal \(\order{1/\varepsilon^2}\) bound for smooth nonconvex optimization, matching the exact-evaluation complexity of regularized Newton methods.
Rather than relying throughout on exact-value acceptance tests, our method uses noisy function values only while they certify a reference decrease and otherwise switches to an objective-function-free regularization scale.
It thereby retains useful line-search behavior without making continued progress depend on unreliable comparisons.

Other closely related methods are OFFO~\citep{Gratton08022024,grattonOFFOMinimizationAlgorithms2023} and Noise Tolerant Quasi-Newton (NTQN) methods~\citep{shiNoiseTolerantQuasiNewtonAlgorithm2022}.
OFFO methods avoid objective-function values altogether, whereas our method retains their use whenever they provide a reliable decrease certificate.
Shi et al.~\citep{shiNoiseTolerantQuasiNewtonAlgorithm2022} develop NTQN for bounded errors in both function and gradient evaluations, using a noise-tolerant line search and a lengthening procedure to stabilize the quasi-Newton updates.
Their analysis establishes linear convergence to a noise-dependent neighborhood for strongly convex objectives while continuing to base step acceptance on noisy function comparisons~\citep[Theorem~3.5]{shiNoiseTolerantQuasiNewtonAlgorithm2022}.
Our relaxed Armijo test is related to their noise-tolerant test~\citep[Section~4]{shiNoiseTolerantQuasiNewtonAlgorithm2022}, but our adaptive error allowance and fallback play a different role.
Function values are used only while they provide a reliable decrease certificate.
This design targets smooth nonconvex problems with accurate gradients and persistent function-value errors, and retains an \(\order{1/\varepsilon^2}\) guarantee for arbitrarily small first-order tolerances.
NTQN instead has the complementary advantage of accommodating bounded gradient errors.
In the function-noise experiments, the proposed variants attain the common gradient tolerance on a larger fraction of instances than NTQN, while the joint-noise results are reported only as stress tests beyond our theory.

The guarantees in \cref{tab:related_work_comparison} have different targets and assumptions.
For line-search L-BFGS with exact evaluations, R-linear convergence of the objective gap holds under uniform convexity and bounded initial scaling matrices~\citep[Theorem~7.1]{liuLimitedMemoryBFGS1989}.
For a regularized Newton method with an Armijo line search, an \(\order{1/\varepsilon^2}\) iteration bound and a uniform bound on the backtracking steps hold under compact-level-set and Lipschitz-Hessian assumptions~\citep[Theorems~4.1 and~4.2]{uedaConvergencePropertiesRegularized2010}.
Kanzow and Steck establish global convergence of regularized L-BFGS to first-order stationarity under mild continuity and boundedness assumptions, but do not give a worst-case rate~\citep[Theorem~2]{kanzowRegularizationLimitedMemory2023}.
By contrast, OFFO and our method have \(\order{1/\varepsilon^2}\) worst-case complexity for first-order stationarity on smooth nonconvex objectives, and our result additionally permits persistent bounded function errors.
\end{MPCchange}

\subsection{Organization of the Paper} % to LLM: ignore this git difference.

\ifSubfilesClassLoaded{
    \bibliographystyle{spbasic}
    \bibliography{../L-BFGS.bib}
}{}

\end{document}
